\section{Discussion}
In HDLSS settings, one effective outlier distance per class summarizes the leading effect of contamination on the hard-margin linear SVM. Under one-sided contamination, distant outliers leave the dual value, weights, and discriminant asymptotically unchanged. Closer outliers cause two distinct failures: below the support-collapse boundary, the contaminated class is represented only by its outliers; below the classification boundary, the intercept moves far enough to misclassify the clean class. Several moderately distant outliers can therefore have the same effect as one close outlier.

Two-sided contamination behaves differently. Outliers on opposite sides pull the limiting intercept in opposite directions. At most one class can collapse and at most one can be misclassified. The boundaries depend on the effective distance in the other class, and symmetric contamination leaves the limiting discriminant at its clean value. The dual optimum still increases and the margin shrinks, so margin diagnostics can reveal contamination even when classification is asymptotically unaffected.

Two practical implications follow. First, the diagonal of the centered Gram matrix estimates each observation's effective distance: it is asymptotically $D_k$ for an outlier and $\Delta/4+\tr(\boldsymbol\Sigma_i)$ for an inlier. Classwise effective distances therefore predict the sign and approximate size of the intercept shift through Theorem~\ref{thm:two-sided}. Second, because much of the damage is mediated by the intercept, a bias correction in the spirit of \citet{nakayama2017}, or a finite soft-margin bound that caps each observation's dual weight, should limit the effect.

Several questions remain. Assumption (A2) requires the outlier offsets to be asymptotically orthogonal to the population contrast. If their inner products are of the same order as $\Delta_*$, cross terms survive and the low-dimensional reductions must be modified. The lower bound in (A4) excludes outliers that lie exactly at the midpoint on the $\Delta_*$ scale; such points dominate the dual and require a different normalization. Extensions to kernel SVMs \citep{nakayama2020}, robust or bias-corrected variants, and comparisons with distance-weighted discrimination \citep{marron2007} are natural directions.
